
\documentclass{article}
\usepackage{colortbl}
\usepackage{makecell}
\usepackage{multirow}
\usepackage{supertabular}

\begin{document}

\newcounter{utterance}

\centering \large Interaction Transcript for game `bbh', experiment `logical\_deduction\_three\_objects', episode 171 with reflexion{-}Qwen3.5{-}27B.
\vspace{24pt}

{ \footnotesize  \setcounter{utterance}{1}
\setlength{\tabcolsep}{0pt}
\begin{supertabular}{c@{$\;$}|p{.15\linewidth}@{}p{.15\linewidth}p{.15\linewidth}p{.15\linewidth}p{.15\linewidth}p{.15\linewidth}}
    \# & \multicolumn{2}{c}{Player} && \multicolumn{2}{c}{Game Master} \\
    \hline

    \theutterance \stepcounter{utterance}  
    & & & \multicolumn{4}{p{0.6\linewidth}}{
        \cellcolor[rgb]{0.9,0.9,0.9}{
            \makecell[{{p{\linewidth}}}]{
                \texttt{\tiny{[P1$\langle$GM]}}
                \texttt{You have attempted this same task before and did not succeed. These are the plans you wrote after those attempts.} \\
\texttt{Trial \#1: New plan: My previous attempt failed because I included a detailed step{-}by{-}step explanation and reasoning before providing the final answer, whereas the few{-}shot examples in the prompt strictly output only the option letter (e.g., "A: (B)"). The task requires deducing the answer internally but presenting only the final selection in the format "A: (Option)". Next time, I will perform the logical deduction silently and output only the final answer string "A: (B)" without any preceding text, explanations, or bullet points.} \\
\\ 
\\ 
\texttt{A logical deduction task which requires deducing the order of a sequence of objects.} \\
\\ 
\texttt{Q: The following paragraphs each describe a set of three objects arranged in a fixed order. The statements are logically consistent within each paragraph. In a golf tournament, there were three golfers: Amy, Eli, and Eve. Eve finished above Amy. Eli finished below Amy.} \\
\texttt{Options:} \\
\texttt{(A) Amy finished last} \\
\texttt{(B) Eli finished last} \\
\texttt{(C) Eve finished last} \\
\texttt{A: (B)} \\
\\ 
\texttt{Q: The following paragraphs each describe a set of three objects arranged in a fixed order. The statements are logically consistent within each paragraph. On a shelf, there are three books: a white book, a green book, and an orange book. The green book is to the right of the white book. The orange book is the rightmost.} \\
\texttt{Options:} \\
\texttt{(A) The white book is the leftmost} \\
\texttt{(B) The green book is the leftmost} \\
\texttt{(C) The orange book is the leftmost} \\
\texttt{A: (A)} \\
\\ 
\texttt{Q: The following paragraphs each describe a set of three objects arranged in a fixed order. The statements are logically consistent within each paragraph. On a shelf, there are three books: a red book, a gray book, and a white book. The white book is to the left of the gray book. The red book is the second from the left.} \\
\texttt{Options:} \\
\texttt{(A) The red book is the leftmost} \\
\texttt{(B) The gray book is the leftmost} \\
\texttt{(C) The white book is the leftmost} \\
\texttt{A: (C)} \\
\\ 
\texttt{Q: The following paragraphs each describe a set of three objects arranged in a fixed order. The statements are logically consistent within each paragraph. On a shelf, there are three books: a yellow book, a brown book, and a red book. The red book is to the left of the brown book. The yellow book is to the right of the brown book.} \\
\texttt{Options:} \\
\texttt{(A) The yellow book is the second from the left} \\
\texttt{(B) The brown book is the second from the left} \\
\texttt{(C) The red book is the second from the left} \\
\texttt{A:} \\
            }
        }
    }
     \\ \\

    \theutterance \stepcounter{utterance}  
    & \multicolumn{4}{p{0.6\linewidth}}{
        \cellcolor[rgb]{0.561,0.737,0.561}{
            \makecell[{{p{\linewidth}}}]{
                \texttt{\tiny{[P1$\rangle$GM]}}
                \texttt{A: (B)} \\
            }
        }
    }
    & & \\ \\

    \theutterance \stepcounter{utterance}  
    & & & \multicolumn{2}{p{0.3\linewidth}}{
        \cellcolor[rgb]{0.9,0.9,0.9}{
            \makecell[{{p{\linewidth}}}]{
                \texttt{\tiny{[GM$|$GM]}}
                \texttt{Target: (B)} \\
            }
        }
    }
    & & \\ \\

    \theutterance \stepcounter{utterance}  
    & & & \multicolumn{2}{p{0.3\linewidth}}{
        \cellcolor[rgb]{0.9,0.9,0.9}{
            \makecell[{{p{\linewidth}}}]{
                \texttt{\tiny{[GM$|$GM]}}
                \texttt{game\_result = LOSE} \\
            }
        }
    }
    & & \\ \\

    \theutterance \stepcounter{utterance}  
    & & \multicolumn{4}{p{0.6\linewidth}}{
        \cellcolor[rgb]{0.95,0.95,0.95}{
            \makecell[{{p{\linewidth}}}]{
                \texttt{\tiny{[P1$\leftrightarrow$P]}}
                \texttt{Post{-}episode Reflexion (not sent to GM):} \\
\texttt{New plan: My previous attempt failed because I incorrectly deduced the order of the books. The prompt states the red book is to the left of the brown book, and the yellow book is to the right of the brown book, establishing the sequence: Red (left), Brown (middle), Yellow (right). Therefore, the brown book is the second from the left, making (B) the correct answer. However, I must verify if my previous output was actually wrong or if the logic was flawed. Wait, in the previous attempt, I output "A: (B)", which corresponds to "The brown book is the second from the left". This is logically correct based on the text: Red < Brown < Yellow. If the outcome was LOSE, it implies the expected answer was different or my internal logic was flawed in a way I didn't catch. Let's re{-}read carefully. "Red is left of Brown" (R < B). "Yellow is right of Brown" (B < Y). Order: R, B, Y. Brown is indeed second. Why did I lose? Perhaps the system expected a different format or I misread the options? Option (B) is "The brown book is the second from the left". This is correct.} \\
\\ 
\texttt{Wait, looking at the "Trial \#1" plan provided in the prompt history, it says: "My previous attempt failed because I included a detailed step{-}by{-}step explanation...". But in the "YOUR LAST ATTEMPT" section, the user output was *only* "A: (B)". This suggests the "Trial \#1" plan refers to a *different* previous attempt not shown in the immediate "YOUR LAST ATTEMPT" block, or the "LOSE" outcome in the current block is due to the *content* being wrong, not the format.} \\
\\ 
\texttt{Let's re{-}evaluate the logic for the current question:} \\
\texttt{1. Red is left of Brown.} \\
\texttt{2. Yellow is right of Brown.} \\
\texttt{Sequence: Red, Brown, Yellow.} \\
\texttt{Positions: 1 (Red), 2 (Brown), 3 (Yellow).} \\
\texttt{Question: Which is the second from the left?} \\
\texttt{Answer: Brown.} \\
\texttt{Option (B): The brown book is the second from the left.} \\
\\ 
\texttt{If I answered (B) and lost, either:} \\
\texttt{1. The logic is actually different (e.g., "right of" means something else? No, standard logic).} \\
\texttt{2. The options are tricky.} \\
\texttt{3. The "LOSE" is because I didn't follow the *exact* format of the few{-}shot examples? The examples show "A: (B)". I output "A: (B)".} \\
\texttt{4. Maybe the previous "Trial \#1" plan is the *only* plan I have, and the current "LOSE" is a result of a different error?} \\
\\ 
\texttt{Actually, looking at the prompt structure, the "Plans from past attempts" section lists} \\
            }
        }
    }
    & \\ \\

\end{supertabular}
}

\end{document}
